\documentclass[12pt]{article}

\usepackage[utf8]{inputenc}
\usepackage{amsmath, amssymb, amsthm}
\usepackage{geometry}
\geometry{a4paper, margin=2.5cm}

\title{Clasă Monotonă}
\author{}
\date{}

\theoremstyle{definition}
\newtheorem{definition}{Definiție}

\begin{document}

\maketitle

\begin{definition}
O familie de mulțimi $\mathcal{M}$ se numește \emph{clasă monotonă} dacă:
\begin{enumerate}
    \item $\Omega \in \mathcal{M}$.
    \item Dacă $A, B \in \mathcal{M}$ și $A \subseteq B$, atunci $B \setminus A \in \mathcal{M}$.
    \item Dacă $(A_n)_{n \ge 1}$ este o familie de mulțimi cu $A_1 \subseteq A_2 \subseteq \cdots$ și $A_n \in \mathcal{M}$ pentru toate $n$, atunci
    

\[
        \bigcup_{n=1}^{\infty} A_n \in \mathcal{M}.
    \]


\end{enumerate}
\end{definition}

\noindent
Aceasta este o caracterizare echivalentă a claselor monotone, utilă în teorema clasei monotone.

\end{document}
